\section{\QSCver{shifted.0.13} Historical Remark Triage and Reader-Hygiene Cleanup}
\label{sec:v401-historical-remark-triage-reader-hygiene-cleanup}

\begin{quote}\small
\noindent\textbf{Status.}
This \QSCver{shifted.0.13} layer is a cleanup and reader-hygiene layer. It does not add a new
theorematic claim, does not close OP~1/4 or OP~5(B), and does not change the
\QSCver{shifted.0.0} global QAM Hilbert-shift. Its purpose is to separate governing shifted
statements from historical genealogy remarks.
\end{quote}

\subsection{Triage principle}
\label{subsec:v401-triage-principle}

\begin{definition}[Historical remark triage classes]
\label{def:v401-historical-remark-triage-classes}
Historical or version-genealogical remarks in the Active Core are assigned one
of the following cleanup roles:
\begin{enumerate}[leftmargin=*]
 \item \textbf{H1: governing guard.} The remark remains in the Active Core
 because it prevents overclaim, preserves OP status, separates Shift A from
 Shift B, or protects tensor/readout grammar.
 \item \textbf{H2: compact genealogy.} The remark may remain only in compressed
 form, usually as a short version-genealogy pointer.
 \item \textbf{H3: witness-only.} The remark belongs to the Historical Witness
 or release genealogy and should not govern the Active Core.
\end{enumerate}
\end{definition}

\begin{definition}[Governing statement criterion]
\label{def:v401-governing-statement-criterion}
A remark is a governing statement only if it changes how an active shifted object is
read, guarded, classified, backtranslated, or status-preserved. A remark that
only explains how a release-final state stage was reached is historical genealogy.
\end{definition}

\begin{remark}[No deletion of proof memory]
\label{rem:v401-no-deletion-proof-memory}
The cleanup does not delete proof memory. It only changes where such memory is
allowed to govern. The Active Core governs through the shifted-namespace normal form; the
Historical Witness preserves genealogy.
\end{remark}

\subsection{Active-Core retention rules}
\label{subsec:v401-active-core-retention-rules}

\begin{proposition}[Active-Core retention rule]
\label{prop:v401-active-core-retention-rule}
A historical remark should remain in the Active Core only if it satisfies at
least one of the following conditions:
\begin{enumerate}[leftmargin=*]
 \item it protects the distinction between local HC-boundary activation
 and global shifted QAM namespace migration;
 \item it protects the distinction between tag, support, packed mask, route
 payload, tensor component, and historical witness value;
 \item it protects tensor grammar, readout target, guard, residual map, or OP
 status;
 \item it points to the Historical Witness as proof memory rather than as a
 second active theorematic line.
\end{enumerate}
\end{proposition}

\begin{proof}
These four cases correspond exactly to the places where historical context
prevents an incorrect active reading. In all other cases, the shifted-namespace normal form
already supplies the governing rule. Therefore long release-final state-stage or
development-stage remarks may be compressed or moved to the witness layer
without changing the mathematics.
\end{proof}

\subsection{Version-genealogy compression}
\label{subsec:v401-version-genealogy-compression}

\begin{definition}[Compressed shifted genealogy note]
\label{def:v401-compressed-shifted-genealogy-note}
The \QSCver{shifted.0.0} migration chain is summarized as follows:
\[
\begin{array}{ll}
$\text{historical release line:}$  & \text{local HC-boundary activation }0_{\rm mask}\to1_{\rm mask},\\
\text{historical release line:} & \text{migration map, backtranslation, and concrete ledger},\\
\text{the local proof line:} & \text{global shifted namespace activation},\\
\text{the local proof line:} & \text{occurrence ledger and shift-audit map},\\
\text{the local proof line:} & \text{pre-shift/shifted/backtranslation display triples},\\
\text{the local proof line:} & \text{physical-readout compatibility check},\\
\text{the local proof line:} & \text{governing shifted-namespace normal form and release-freeze handoff}.
\end{array}
\]
This note replaces repeated long-form release-final state-stage narrative in the
governing reading of the Active Core.
\end{definition}

\begin{remark}[Current governing reading]
\label{rem:v401-current-governing-reading}
In the Active Core, shifted is no longer a future release-final state language. It is the
current governing namespace for active QAM/router/readout objects. Historical
and pre-shift material remains readable through the Historical Witness and the
backtranslation ledger.
\end{remark}

\subsection{Cleanup audit result}
\label{subsec:v401-cleanup-audit-result}

\begin{theorem}[Historical remark cleanup conservativity]
\label{thm:v401-historical-remark-cleanup-conservativity}
The \QSCver{shifted.0.13} historical remark triage is conservative over \QSCver{shifted.0.0}: it does not
change the shifted namespace, the backtranslation map, the occurrence classes,
the display triples, the physical-readout compatibility checks, or the
governing shifted-namespace normal form.
\end{theorem}

\begin{proof}
The cleanup assigns roles to historical remarks and compresses genealogy into a
single note. It does not modify the definitions of
\[
 \Sigma_4(k)=k+1,\qquad
 \mathrm B_4(a)=a-1,\qquad
 \mathcal N_4(X)=(C_X,X_{\mathrm{shift}},\mathrm B_4(X_{\mathrm{shift}}),P_X,G_X,S_X).
\]
It also leaves the non-claim guards for OP~1/4 and OP~5(B) unchanged. Hence
the mathematical content of \QSCver{shifted.0.0} is preserved.
\end{proof}

\begin{remark}[\QSCver{shifted.0.13} non-claim]
\label{rem:v401-non-claim}
Version~4.0.13 is a release-final cleanup version. It does not add new physics, does not
close OP~1/4 or OP~5(B), does not activate primitive mixed tensors, and does
not claim Turing-completeness.
\end{remark}



\clearpage
\section{\QSCver{shifted.0.13} Reader Navigation Cleanup}
\label{sec:v402-reader-navigation-cleanup}

\begin{quote}\small
\noindent\textbf{Status.}
This \QSCver{shifted.0.13} layer is a reader-navigation cleanup. It does not change the shifted
shift, the occurrence ledger, display triples, readout compatibility, normal
form, OP status, or Historical Witness.
\end{quote}

\subsection{Navigation principle}
\label{subsec:v402-navigation-principle}

\begin{definition}[Current reader-navigation principle]
\label{def:v402-current-reader-navigation-principle}
The current reader-navigation principle is:
\[
 \text{Active Core governs},\qquad
 \text{Historical Witness preserves},\qquad
 \text{backtranslation connects}.
\]
Thus active shifted statements are read in the shifted-namespace normal form, historical material is
read as proof memory, and pre-shift displays are connected to active shifted displays
only through explicit backtranslation records.
\end{definition}

\begin{definition}[Governing-map roles]
\label{def:v402-governing-map-roles}
The following roles are used for reader navigation:
\begin{enumerate}[leftmargin=*]
 \item \textbf{Governing active role:} current shifted definitions, theorems,
 normal forms, guards, and OP-status statements.
 \item \textbf{Routing role:} tables, display triples, occurrence classes, and
 backtranslation ledgers that explain how to read active objects.
 \item \textbf{Genealogy role:} historical remarks and witness material that
 explain how the current line was reached.
\end{enumerate}
\end{definition}

\begin{lemma}[Navigation cleanup is conservative]
\label{lem:v402-navigation-cleanup-conservative}
The \QSCver{shifted.0.13} reader-navigation cleanup is conservative over \QSCver{shifted.0.1}.
\end{lemma}

\begin{proof}
The cleanup adds navigation instructions and clarifies the current governing
reading. It does not modify
\[
 \Sigma_4(k)=k+1,\qquad
 \mathrm B_4(a)=a-1,\qquad
 \mathcal N_4(X)=(C_X,X_{\mathrm{shift}},\mathrm B_4(X_{\mathrm{shift}}),P_X,G_X,S_X).
\]
It also preserves the Historical Witness and leaves all OP-status guards
unchanged. Therefore it is conservative over \QSCver{shifted.0.1}.
\end{proof}

\begin{remark}[\QSCver{shifted.0.13} non-claim]
\label{rem:v402-non-claim}
Version~4.0.13 is a navigation release-final cleanup version. It does not add new physics,
does not close OP~1/4 or OP~5(B), does not activate primitive mixed tensors,
and does not claim Turing-completeness.
\end{remark}



\clearpage
\section{\QSCver{shifted.0.13} Notation and Symbol Hygiene Cleanup}
\label{sec:v403-notation-symbol-hygiene-cleanup}

\begin{quote}\small
\noindent\textbf{Status.}
This \QSCver{shifted.0.13} layer is a notation-hygiene cleanup. It does not change the shifted
namespace shift, the backtranslation ledger, the occurrence classes, the tensor
grammar, the readout compatibility checks, OP status, or the Historical
Witness. It records canonical symbol usage for the current shifted line.
\end{quote}

\subsection{Canonical shifted symbol registry}
\label{subsec:v403-canonical-shifted-symbol-registry}

\begin{definition}[Canonical shifted symbol registry]
\label{def:v403-canonical-shifted-symbol-registry}
The current shifted Active Core uses the following canonical symbols:
\[
\begin{array}{lll}
\Sigma_4(k)=k+1
 & \text{global shift on inhabited pre-shift tags}
 & \text{language/router namespace},\\[0.2em]
\mathrm B_4(a)=a-1
 & \text{backtranslation on inhabited shifted tags }a\geq 1
 & \text{partial inverse},\\[0.2em]
0_{\mathrm{shift}}=\varnothing_{\mathrm{active}}
 & \text{reserved empty/no-active-kind marker}
 & \text{not inhabited},\\[0.2em]
0_{\mathrm{mask}}\to1_{\mathrm{mask}}
 & \text{local HC-boundary activation}
 & \text{Shift A},\\[0.2em]
\mathfrak D_4(X)
 & \text{pre-shift/shifted/backtranslation display triple}
 & \text{reader/audit display},\\[0.2em]
\mathcal N_4(X)
 & \text{governing shifted normal-form object}
 & \text{active object form}.
\end{array}
\]
\end{definition}

\begin{remark}[Roman backtranslation symbol]
\label{rem:v403-roman-backtranslation-symbol}
The backtranslation symbol is written as $\mathrm B_4$, not as a Greek beta
symbol. This avoids confusion with an additional operator, physical beta
parameter, or unstated tensor component.
\end{remark}

\subsection{Disjoint role vocabulary}
\label{subsec:v403-disjoint-role-vocabulary}

\begin{definition}[Tag, route, support, packed mask, and tensor payload]
\label{def:v403-role-vocabulary}
The following terms are kept disjoint in the shifted Active Core:
\begin{enumerate}[leftmargin=*]
 \item A \emph{tag} is an inhabited QAM/router label subject to the shifted
 namespace when active.
 \item A \emph{route payload} is an ordered finite list of tags or tag-like
 route entries.
 \item A \emph{support set} is an unordered finite set of inhabited support
 entries.
 \item A \emph{packed mask integer} is the integer encoding of a support set
 and is not itself shifted as a tag.
 \item A \emph{tensor-index payload} is a QAM/router payload carried by a
 two-sided tensor index position.
 \item A \emph{tensor component} is not automatically a QAM/router tag and is
 shifted only if explicitly declared as a tag payload.
\end{enumerate}
\end{definition}

\begin{lemma}[Packed mask is not a tag]
\label{lem:v403-packed-mask-is-not-tag}
A packed mask integer is not shifted by $\Sigma_4$ as an integer.
\end{lemma}

\begin{proof}
By Definition~\ref{def:v403-role-vocabulary}, a packed mask integer encodes a
support set. The shifted shift acts on inhabited support entries. The packed
integer is then recomputed from the shifted support. Therefore the operation is
not $m\mapsto m+1$ on the packed integer.
\end{proof}

\subsection{Canonical display forms}
\label{subsec:v403-canonical-display-forms}

\begin{definition}[Canonical display forms]
\label{def:v403-canonical-display-forms}
The following display forms are canonical for current shifted cleanup work:
\[
 \mathfrak D_4(X)
 =
 (X_{\mathrm{pre}},X_{\mathrm{shift}},\mathrm B_4(X_{\mathrm{shift}})),
\]
and
\[
 \mathcal N_4(X)
 =
 (C_X,X_{\mathrm{shift}},\mathrm B_4(X_{\mathrm{shift}}),P_X,G_X,S_X).
\]
Here $C_X$ is the occurrence class, $P_X$ is the payload or readout target,
$G_X$ is the guard record, and $S_X$ is the theorematic status record.
\end{definition}

\begin{remark}[Display triple versus normal form]
\label{rem:v403-display-triple-versus-normal-form}
The display triple $\mathfrak D_4(X)$ is a reader/audit surface. The normal
form $\mathcal N_4(X)$ is the governing active object form. They should not be
used interchangeably.
\end{remark}

\subsection{Tensor notation hygiene, declared local OP1/OP4 closure, release-final}
\label{subsec:v403-tensor-notation-hygiene}

\begin{definition}[Tensor payload hygiene rule]
\label{def:v403-tensor-payload-hygiene-rule}
The shifted shift may act on QAM/router payloads carried inside tensor-index
positions. It does not act on sidedness, upper/lower variance, ordinary
upper-lower contraction, or explicit same-variance comparison rules.
\end{definition}

\begin{proposition}[Tensor hygiene preservation]
\label{prop:v403-tensor-hygiene-preservation}
The notation-hygiene cleanup preserves the two-sided structural tensor grammar.
\end{proposition}

\begin{proof}
The cleanup only records the distinction between tag payload, tensor-index
position, and tensor component. It does not move indices, change variance, or
alter contraction/comparison rules. Therefore the tensor grammar is preserved.
\end{proof}

\subsection{Notation-hygiene conservativity}
\label{subsec:v403-notation-hygiene-conservativity}

\begin{theorem}[Notation-hygiene conservativity]
\label{thm:v403-notation-hygiene-conservativity}
The \QSCver{shifted.0.13} notation-hygiene cleanup is conservative over \QSCver{shifted.0.2}: it changes no
mathematical content of the global shift, backtranslation, occurrence ledger,
display triples, readout compatibility, normal form, or OP-status guards.
\end{theorem}

\begin{proof}
The cleanup introduces a registry and role vocabulary for existing objects.
It leaves the definitions of $\Sigma_4$, $\mathrm B_4$, $\mathfrak D_4$, and
$\mathcal N_4$ unchanged. It also preserves the distinction between
$0_{\mathrm{shift}}$ and $0_{\mathrm{mask}}$, and it does not modify OP-status
guards. Hence the cleanup is conservative.
\end{proof}

\begin{remark}[\QSCver{shifted.0.13} non-claim]
\label{rem:v403-non-claim}
Version~4.0.13 is a notation release-final cleanup version. It does not add new physics,
does not close OP~1/4 or OP~5(B), does not activate primitive mixed tensors,
and does not claim Turing-completeness.
\end{remark}



\clearpage
\section{\QSCver{shifted.0.13} Table, Display, and Example Normalization Cleanup}
\label{sec:v404-table-display-example-normalization-cleanup}

\begin{quote}\small
\noindent\textbf{Status.}
This \QSCver{shifted.0.13} layer is a table/display/example normalization cleanup. It does
not change the shifted namespace shift, the backtranslation map, the occurrence
classes, the tensor grammar, readout compatibility, OP status, or the
Historical Witness. It standardizes how examples should be displayed.
\end{quote}

\subsection{Canonical example display contract}
\label{subsec:v404-canonical-example-display-contract}

\begin{definition}[Canonical example display contract]
\label{def:v404-canonical-example-display-contract}
Every active shifted example that compares an old display with a current display
should use the three-column contract
\[
 \vfourTableHead .
\]
Equivalently, the example should be readable as the display triple
\[
 \mathfrak D_4(X)
 =
 \vfourDisplayTriple
 {X_{\mathrm{pre}}}
 {X_{\mathrm{shift}}}
 {\mathrm B_4(X_{\mathrm{shift}})}.
\]
\end{definition}

\begin{remark}[Display contract is not a theorem]
\label{rem:v404-display-contract-not-theorem}
The display contract is a reader/audit convention. It does not add a proof
certificate and does not change theorematic status.
\end{remark}

\subsection{Normalized route and support examples}
\label{subsec:v404-normalized-route-support-examples}

\begin{center}
\small
\setlength{\tabcolsep}{4pt}
\begin{tabular}{@{}lll@{}}
\toprule
pre-shift source & shifted display & backtranslation\\
\midrule
tag $5$ & tag $6$ & $\mathrm B_4(6)=5$\\
tag $12$ & tag $13$ & $\mathrm B_4(13)=12$\\
route $\langle5,12\rangle$ & route $\langle6,13\rangle$ & route $\langle5,12\rangle$\\
support $\{5,12\}$ & support $\{6,13\}$ & support $\{5,12\}$\\
\bottomrule
\end{tabular}
\end{center}

\begin{lemma}[Route/support display normalization]
\label{lem:v404-route-support-display-normalization}
The normalized route and support examples preserve the shifted backtranslation
ledger on inhabited entries.
\end{lemma}

\begin{proof}
Each shifted entry in the displayed route or support is at least $1$ and therefore
has backtranslation $\mathrm B_4(a)=a-1$. Applying this componentwise returns
the pre-shift route or support.
\end{proof}

\subsection{Normalized packed-mask example}
\label{subsec:v404-normalized-packed-mask-example}

\begin{center}
\small
\setlength{\tabcolsep}{4pt}
\begin{tabular}{@{}lll@{}}
\toprule
pre-shift source & shifted display & backtranslation\\
\midrule
support $S=\{5,12\}$ & support $S^{(4)}=\{6,13\}$ & support $S=\{5,12\}$\\
packed mask $m_3(S)=4128$ & packed mask $m_4(S^{(4)})=8256$ & via support, not by integer subtraction\\
\bottomrule
\end{tabular}
\end{center}

\begin{proposition}[Packed-mask display normalization]
\label{prop:v404-packed-mask-display-normalization}
A packed-mask display is normalized only when it separates support migration
from packed-integer recomputation.
\end{proposition}

\begin{proof}
The shifted shift acts on inhabited support entries, not on the packed integer as a
tag. Once the support $S$ is shifted to $S^{(4)}$, the packed integer is
recomputed:
\[
 m_4(S^{(4)})=\sum_{k\in S}2^{k+1}=2\sum_{k\in S}2^k=2m_3(S).
\]
Thus the normalized display must show support migration separately from the
packed integer.
\end{proof}

\subsection{Normalized tensor-display examples}
\label{subsec:v404-normalized-tensor-display-examples}

\begin{center}
\scriptsize
\setlength{\tabcolsep}{2pt}
\begin{tabular}{@{}p{0.30\textwidth}P{0.30\textwidth}P{0.30\textwidth}@{}}
\toprule
pre-shift source & shifted display & backtranslation\\
\midrule
${}^{\langle5,12\rangle}{}_{\langle2\rangle}\mathcal T^{i}{}_{j}$
&
${}^{\langle6,13\rangle}{}_{\langle3\rangle}\mathcal T^{i}{}_{j}$
&
${}^{\langle5,12\rangle}{}_{\langle2\rangle}\mathcal T^{i}{}_{j}$
\\[0.5em]
${}^{\langle5,12\rangle}{}_{\langle2\rangle}\mathcal T^{\langle7\rangle}{}_{\langle10\rangle}$
&
${}^{\langle6,13\rangle}{}_{\langle3\rangle}\mathcal T^{\langle8\rangle}{}_{\langle11\rangle}$
&
${}^{\langle5,12\rangle}{}_{\langle2\rangle}\mathcal T^{\langle7\rangle}{}_{\langle10\rangle}$
\\
\bottomrule
\end{tabular}
\end{center}

\begin{lemma}[Tensor-display normalization preserves grammar]
\label{lem:v404-tensor-display-normalization-preserves-grammar}
The normalized tensor-display examples preserve sidedness, variance,
contraction rules, and same-variance comparison rules.
\end{lemma}

\begin{proof}
Only QAM/router payloads inside index positions are shifted. No index is moved
between front/back positions or between upper/lower variance. Therefore the
two-sided tensor grammar is unchanged.
\end{proof}

\subsection{Example normalization rule}
\label{subsec:v404-example-normalization-rule}

\begin{definition}[Example normalization rule]
\label{def:v404-example-normalization-rule}
A shifted example is normalized if it satisfies all of the following:
\begin{enumerate}[leftmargin=*]
 \item it declares whether the displayed object is a tag, route payload,
 support set, packed mask integer, tensor-index payload, tensor component, or
 historical witness value;
 \item it displays the shifted form and the backtranslation when a shifted active
 tag payload is present;
 \item it separates support migration from packed-integer recomputation;
 \item it preserves tensor sidedness and variance in tensor examples;
 \item it preserves guard and status information in OP-facing examples.
\end{enumerate}
\end{definition}

\begin{theorem}[Example-normalization conservativity]
\label{thm:v404-example-normalization-conservativity}
The \QSCver{shifted.0.13} table/display/example normalization cleanup is conservative over
\QSCver{shifted.0.3}.
\end{theorem}

\begin{proof}
The cleanup standardizes how examples are shown. It does not change the
definitions of $\Sigma_4$, $\mathrm B_4$, $\mathfrak D_4$, or
$\mathcal N_4$. It also leaves tensor grammar, readout guards, OP status, and
the Historical Witness unchanged. Therefore it is conservative over \QSCver{shifted.0.3}.
\end{proof}

\begin{remark}[\QSCver{shifted.0.13} non-claim]
\label{rem:v404-non-claim}
Version~4.0.13 is a table/display/example normalization release-final state. It does not
add new physics, does not close OP~1/4 or OP~5(B), does not activate primitive
mixed tensors, and does not claim Turing-completeness.
\end{remark}



\clearpage
\section{\QSCver{shifted.0.13} Appendix and Historical Witness Linkage Cleanup}
\label{sec:v405-appendix-historical-witness-linkage-cleanup}

\begin{quote}\small
\noindent\textbf{Status.}
This \QSCver{shifted.0.13} layer is an appendix and Historical Witness linkage cleanup. It
does not change the shifted namespace shift, the backtranslation map, the occurrence
classes, display triples, tensor grammar, readout compatibility, OP status, or
the Historical Witness. It clarifies how the Active Core and Historical
Witness should be read together.
\end{quote}

\subsection{Split-bundle linkage principle}
\label{subsec:v405-split-bundle-linkage-principle}

\begin{definition}[Split-bundle linkage principle]
\label{def:v405-split-bundle-linkage-principle}
The split bundle is read by the principle
\[
 \text{Active Core governs},\qquad
 \text{Historical Witness preserves},\qquad
 \text{backtranslation connects}.
\]
The Active Core contains the current governing shifted language and theorematic
status. The Historical Witness preserves proof memory, genealogy, and older
readout surfaces. The backtranslation ledger explains how inhabited shifted
shifted tags relate to pre-shift tags.
\end{definition}

\begin{remark}[No second active theory line]
\label{rem:v405-no-second-active-theory-line}
The Historical Witness is not a second active theory line. If a historical
statement conflicts with the current shifted-namespace normal form, the Active Core governs
the current reading, while the Historical Witness records the genealogy of the
older formulation.
\end{remark}

\subsection{Historical Witness use classes}
\label{subsec:v405-historical-witness-use-classes}

\begin{definition}[Historical Witness use classes]
\label{def:v405-historical-witness-use-classes}
Historical Witness material is used in one of the following roles:
\begin{enumerate}[leftmargin=*]
 \item \textbf{Proof-memory role.} The material preserves derivation history,
 older proof attempts, and reconstruction context.
 \item \textbf{Genealogy role.} The material explains why a current shifted object
 has the form it has.
 \item \textbf{Backtranslation role.} The material provides a pre-shift source
 display that can be related to a current shifted display.
 \item \textbf{Non-governing archive role.} The material is retained for
 auditability but does not govern current active statements.
\end{enumerate}
\end{definition}

\begin{lemma}[Witness roles preserve active governance]
\label{lem:v405-witness-roles-preserve-active-governance}
The Historical Witness use classes preserve the Active Core as the governing
current line.
\end{lemma}

\begin{proof}
Each use class assigns historical material a support role: proof memory,
genealogy, backtranslation, or archive. None of these roles overrides the
current shifted-namespace normal form, occurrence ledger, readout guard, or OP-status record.
Therefore active governance remains in the Active Core.
\end{proof}

\subsection{Crosswalk table}
\label{subsec:v405-crosswalk-table}

\begin{center}
\scriptsize
\setlength{\tabcolsep}{2pt}
\begin{tabular}{@{}p{0.20\textwidth}P{0.31\textwidth}P{0.34\textwidth}@{}}
\toprule
Reader question & Active Core answer & Historical Witness use\\
\midrule
What is current? & shifted-namespace normal form and current guards & not governing\\
Where did it come from? & compressed genealogy note & proof-memory/genealogy source\\
How do old tags read? & $\mathrm B_4(a)=a-1$ on inhabited shifted tags & pre-shift display source\\
What is the OP status? & current OP-status guard & older attempts remain context\\
What about older examples? & read only if reintroduced by display triple & archive/backtranslation source\\
\bottomrule
\end{tabular}
\end{center}

\subsection{Appendix linkage rule}
\label{subsec:v405-appendix-linkage-rule}

\begin{definition}[Appendix linkage rule]
\label{def:v405-appendix-linkage-rule}
An appendix or historical block may be linked to the Active Core only by one of
the following explicit link types:
\[
 \text{genealogy link},\qquad
 \text{backtranslation link},\qquad
 \text{proof-memory link},\qquad
 \text{audit link}.
\]
A link is not a promotion unless the Active Core explicitly supplies the
corresponding shifted normal-form object and status guard.
\end{definition}

\begin{proposition}[Appendix links are not promotions]
\label{prop:v405-appendix-links-are-not-promotions}
An appendix or Historical Witness link does not by itself promote historical
material to current active theorematic status.
\end{proposition}

\begin{proof}
By Definition~\ref{def:v405-appendix-linkage-rule}, a link has one of four
support roles. None of those roles is a promotion certificate. Promotion to
current active status requires an explicit shifted normal-form object and a status
guard in the Active Core. Therefore the link alone does not promote historical
material.
\end{proof}

\subsection{Historical tag and example reading}
\label{subsec:v405-historical-tag-example-reading}

\begin{definition}[Historical tag reading]
\label{def:v405-historical-tag-reading}
A tag, route, support, packed mask, or tensor-payload example inside the
Historical Witness is not automatically a current shifted object. It becomes
current only if the Active Core reintroduces it through
\[
 \mathfrak D_4(X)
 =
 (X_{\mathrm{pre}},X_{\mathrm{shift}},\mathrm B_4(X_{\mathrm{shift}}))
\]
or through the governing normal form
\[
 \mathcal N_4(X)
 =
 (C_X,X_{\mathrm{shift}},\mathrm B_4(X_{\mathrm{shift}}),P_X,G_X,S_X).
\]
\end{definition}

\begin{remark}[Historical examples are protected from silent shift]
\label{rem:v405-historical-examples-protected-from-silent-shift}
Historical examples are not silently shifted. They are either preserved as
witness material or explicitly connected to the current shifted line by a display
triple or normal-form record.
\end{remark}

\subsection{Historical linkage conservativity}
\label{subsec:v405-historical-linkage-conservativity}

\begin{theorem}[Historical Witness linkage conservativity]
\label{thm:v405-historical-witness-linkage-conservativity}
The \QSCver{shifted.0.13} Historical Witness linkage cleanup is conservative over \QSCver{shifted.0.4}: it
does not change the current shifted namespace, backtranslation ledger, occurrence
classes, display contracts, tensor grammar, readout compatibility, OP status,
or Historical Witness content.
\end{theorem}

\begin{proof}
The cleanup adds link roles and reading rules between the Active Core and
Historical Witness. It does not modify the definitions of $\Sigma_4$,
$\mathrm B_4$, $\mathfrak D_4$, or $\mathcal N_4$. It does not alter any
theorematic status record or OP guard. It also does not remove Historical
Witness material. Hence the cleanup is conservative over \QSCver{shifted.0.4}.
\end{proof}

\begin{remark}[\QSCver{shifted.0.13} non-claim]
\label{rem:v405-non-claim}
Version~4.0.13 is an appendix and Historical Witness linkage release-final cleanup version.
It does not add new physics, does not close OP~1/4 or OP~5(B), does not
activate primitive mixed tensors, and does not claim Turing-completeness.
\end{remark}



\clearpage
\section{\QSCver{shifted.0.13} Clean Reader Bundle and Hygiene-Freeze Candidate}
\label{sec:v406-clean-reader-bundle-hygiene-freeze-release-final state}

\begin{quote}\small
\noindent\textbf{Status.}
This \QSCver{shifted.0.13} layer closes the local cleanup sequence as a clean reader
bundle release-final state. It does not add a new proof claim, does not close OP~1/4 or
OP~5(B), does not activate primitive mixed tensors, and does not change the
\QSCver{shifted.0.0} global QAM Hilbert-shift. It records that the cleanup line from
\QSCver{shifted.0.1} to \QSCver{shifted.0.5} has stabilized the reader-facing structure.
\end{quote}

\subsection{Clean reader bundle principle}
\label{subsec:v406-clean-reader-bundle-principle}

\begin{definition}[Clean reader bundle principle]
\label{def:v406-clean-reader-bundle-principle}
A clean shifted reader bundle is a split bundle in which:
\begin{enumerate}[leftmargin=*]
 \item the Active Core is the current governing line;
 \item the Historical Witness is preserved as proof memory and genealogy;
 \item active QAM/router/readout objects are read in shifted-namespace normal form;
 \item historical/pre-shift material is connected only by explicit witness,
 genealogy, audit, or backtranslation links;
 \item examples and tables use normalized display contracts whenever they
 compare pre-shift and shifted displays.
\end{enumerate}
\end{definition}

\begin{remark}[Clean bundle is not a shorter theory]
\label{rem:v406-clean-bundle-not-shorter-theory}
The clean reader bundle is not a shortened or weakened theory. It is the same
shifted governing structure with improved reader navigation, symbol hygiene,
example hygiene, and Historical Witness linkage.
\end{remark}

\subsection{cleanup line cleanup closure map}
\label{subsec:v406-cleanup-closure-map}

\begin{center}
\scriptsize
\setlength{\tabcolsep}{2pt}
\begin{tabular}{@{}p{0.16\textwidth}P{0.36\textwidth}P{0.36\textwidth}@{}}
\toprule
Version & Cleanup role & Mathematical status\\
\midrule
\QSCver{shifted.0.1} & historical remark triage & conservative; no theorematic change\\
\QSCver{shifted.0.2} & reader navigation map & conservative; current reading clarified\\
\QSCver{shifted.0.3} & notation and symbol hygiene & conservative; symbol roles registered\\
\QSCver{shifted.0.4} & table/display/example normalization & conservative; examples standardized\\
\QSCver{shifted.0.5} & Historical Witness linkage & conservative; split-bundle roles clarified\\
\QSCver{shifted.0.13} & clean reader bundle release-final state & conservative; hygiene line frozen locally\\
\bottomrule
\end{tabular}
\end{center}

\subsection{Reader-facing invariant checklist}
\label{subsec:v406-reader-facing-invariant-checklist}

\begin{definition}[Reader-facing invariant checklist]
\label{def:v406-reader-facing-invariant-checklist}
The following invariants are preserved by the cleanup sequence:
\[
\begin{array}{ll}
\text{shift:} & \Sigma_4(k)=k+1,\\[0.2em]
\text{backtranslation:} & \mathrm B_4(a)=a-1\quad(a\geq1),\\[0.2em]
\text{reserved zero:} & 0_{\mathrm{shift}}=\varnothing_{\mathrm{active}},\\[0.2em]
$\text{HC-boundary shift:}$  & 0_{\mathrm{mask}}\to1_{\mathrm{mask}},\\[0.2em]
\text{display triple:} & \mathfrak D_4(X)=(X_{\mathrm{pre}},X_{\mathrm{shift}},\mathrm B_4(X_{\mathrm{shift}})),\\[0.2em]
\text{normal form:} & \mathcal N_4(X)=(C_X,X_{\mathrm{shift}},\mathrm B_4(X_{\mathrm{shift}}),P_X,G_X,S_X).
\end{array}
\]
\end{definition}

\begin{lemma}[Cleanup sequence preserves the reader-facing invariants]
\label{lem:v406-cleanup-sequence-preserves-reader-facing-invariants}
The \QSCver{shifted.0.1}--\QSCver{shifted.0.13} cleanup sequence preserves every invariant listed in
Definition~\ref{def:v406-reader-facing-invariant-checklist}.
\end{lemma}

\begin{proof}
The cleanup sequence adds reader navigation, historical remark triage, notation
hygiene, example normalization, and Historical Witness linkage. None of these
operations alters the formulas in the invariant checklist. Therefore the
reader-facing invariants are preserved.
\end{proof}

\subsection{Clean reader bundle handoff}
\label{subsec:v406-clean-reader-bundle-handoff}

\begin{definition}[Clean reader bundle handoff]
\label{def:v406-clean-reader-bundle-handoff}
After \QSCver{shifted.0.13}, new work should be separated into one of the following:
\begin{enumerate}[leftmargin=*]
 \item \textbf{further cleanup:} layout, bibliography, index, navigation, or
 display polishing without mathematical change;
 \item \textbf{active shifted development:} new QAM/router/readout objects written
 directly in the shifted-namespace normal form;
 \item \textbf{release packaging:} DOI insertion, final metadata, and Zenodo
 release notes.
\end{enumerate}
\end{definition}

\begin{remark}[Preferred next mathematical branch]
\label{rem:v406-preferred-next-mathematical-branch}
The preferred next mathematical branch after a clean reader bundle is not more
namespace work. It is active use of the shifted-namespace normal form on OP-facing readout
packages, while preserving the OP-status guards.
\end{remark}

\subsection{Clean reader bundle conservativity}
\label{subsec:v406-clean-reader-bundle-conservativity}

\begin{theorem}[Clean reader bundle conservativity]
\label{thm:v406-clean-reader-bundle-conservativity}
The \QSCver{shifted.0.13} historical source-hygiene layer is conservative over the \QSCver{shifted.0.0} release
final and over the \QSCver{shifted.0.1}--\QSCver{shifted.0.5} cleanup sequence.
\end{theorem}

\begin{proof}
The release-final state records the accumulated cleanup discipline without changing the
shifted shift, backtranslation ledger, occurrence classes, display triples, tensor
grammar, readout compatibility checks, Historical Witness content, or OP-status
guards. It only summarizes and freezes the reader-facing cleanup structure.
Therefore it is conservative over \QSCver{shifted.0.0} and over the intervening cleanup
sequence.
\end{proof}

\begin{remark}[\QSCver{shifted.0.13} non-claim]
\label{rem:v406-non-claim}
Version~4.0.13 is a clean reader bundle release-final state. It does not add new physics,
does not close OP~1/4 or OP~5(B), does not activate primitive mixed tensors,
and does not claim Turing-completeness.
\end{remark}



\clearpage
\section{\QSCver{shifted.0.13} Frontmatter, Metadata, and DOI-Record Hygiene Cleanup}
\label{sec:v407-frontmatter-metadata-doi-record-hygiene-cleanup}

\begin{quote}\small
\noindent\textbf{Status.}
This \QSCver{shifted.0.13} layer is a frontmatter and metadata hygiene cleanup. It does not
change the shifted namespace shift, the backtranslation map, the occurrence classes,
display triples, tensor grammar, readout compatibility, OP status, or the
Historical Witness. It records a stable rule for title, DOI, version, and
release-final state/release wording.
\end{quote}

\subsection{Frontmatter invariants}
\label{subsec:v407-frontmatter-invariants}

\begin{definition}[Frontmatter invariants]
\label{def:v407-frontmatter-invariants}
The following frontmatter invariants are fixed for the cleanup-line:
\begin{enumerate}[leftmargin=*]
 \item the title remains
 \[
 \textit{Quantum Sphaera Companion: A Structure-First Mathematical Unfolding};
 \]
 \item the Concept DOI remains the concept-level DOI;
 \item the public \QSCver{shifted.0.0} Version DOI remains a prior version DOI for local
 release-final cleanup versions unless a new Zenodo DOI is explicitly assigned;
 \item local release-final cleanup versions must not present themselves as release-final
 DOI-bearing versions;
 \item the Historical Witness remains separately preserved in the bundle.
\end{enumerate}
\end{definition}

\begin{remark}[Candidate wording]
\label{rem:v407-release-final state-wording}
The word ``release-final state'' is acceptable in local cleanup line cleanup metadata and file
names. It should not appear as a mathematical status operator and should not
be used to weaken already governing shifted definitions.
\end{remark}

\subsection{DOI-record hygiene}
\label{subsec:v407-doi-record-hygiene}

\begin{definition}[DOI-record hygiene rule]
\label{def:v407-doi-record-hygiene-rule}
For this release-final cleanup version, the DOI record is read as
\[
\begin{array}{ll}
\text{Concept DOI:} & 10.5281/zenodo.19210728,\\
\text{Prior public version DOI:} & 10.5281/zenodo.19805670,\\
\text{Version DOI:} & 10.5281/zenodo.19811662.
\end{array}
\]
The version DOI was inserted during release-final packaging.
\end{definition}

\begin{lemma}[DOI hygiene preserves release genealogy]
\label{lem:v407-doi-hygiene-preserves-release-genealogy}
The DOI-record hygiene rule preserves the public \QSCver{shifted.0.0} release genealogy while
allowing local release-final cleanup versions to remain unassigned.
\end{lemma}

\begin{proof}
The concept DOI points to the concept-level record. The prior public version
DOI points to the last DOI-bearing release. Since a local release-final cleanup version has
no assigned assigned version DOI, no new DOI-bearing release is asserted. Therefore
release genealogy remains correct.
\end{proof}

\subsection{Metadata role separation}
\label{subsec:v407-metadata-role-separation}

\begin{definition}[Metadata role separation]
\label{def:v407-metadata-role-separation}
Version labels, file names, DOI records, and release notes are metadata. They
do not define mathematical objects. Mathematical objects must be named by their
mathematical roles, such as $\Sigma_4$, $\mathrm B_4$, $\mathfrak D_4$, or
$\mathcal N_4$, not by release-file names.
\end{definition}

\begin{proposition}[Metadata does not alter theorematic status]
\label{prop:v407-metadata-does-not-alter-theorematic-status}
Changing a local cleanup version label from historical release marker \QSCver{shifted.0.6} to \QSCver{shifted.0.13} does not alter
the theorematic status of any shifted object.
\end{proposition}

\begin{proof}
The version label is metadata. The theorematic status of a shifted object is stored
in its governing status record, for example the $S_X$ field in
$\mathcal N_4(X)$. Updating a cleanup label does not modify that status
record. Hence theorematic status is unchanged.
\end{proof}

\subsection{Frontmatter hygiene checklist}
\label{subsec:v407-frontmatter-hygiene-checklist}

\begin{definition}[Frontmatter hygiene checklist]
\label{def:v407-frontmatter-hygiene-checklist}
A clean local release-final cleanup version should satisfy:
\begin{enumerate}[leftmargin=*]
 \item unchanged title;
 \item visible cleanup version label;
 \item Concept DOI displayed as concept pointer;
 \item prior public version DOI displayed as prior DOI;
 \item no unassigned DOI represented as current DOI;
 \item Active Core and Historical Witness both present;
 \item non-claim guards preserved.
\end{enumerate}
\end{definition}

\subsection{Frontmatter metadata conservativity}
\label{subsec:v407-frontmatter-metadata-conservativity}

\begin{theorem}[Frontmatter metadata conservativity]
\label{thm:v407-frontmatter-metadata-conservativity}
The \QSCver{shifted.0.13} frontmatter/metadata hygiene cleanup is conservative over \QSCver{shifted.0.6}.
\end{theorem}

\begin{proof}
The cleanup records metadata rules and DOI-record hygiene. It does not modify
the definitions of $\Sigma_4$, $\mathrm B_4$, $\mathfrak D_4$, or
$\mathcal N_4$. It does not alter tensor grammar, readout compatibility,
Historical Witness content, or OP-status guards. Therefore it is conservative
over \QSCver{shifted.0.6}.
\end{proof}

\begin{remark}[\QSCver{shifted.0.13} non-claim]
\label{rem:v407-non-claim}
Version~4.0.13 is a frontmatter and metadata hygiene release-final state. It does not add
new physics, does not close OP~1/4 or OP~5(B), does not activate primitive
mixed tensors, and does not claim Turing-completeness.
\end{remark}



\clearpage
\section{\QSCver{shifted.0.13} Bibliography, Citation, and Terminal Record Hygiene Cleanup}
\label{sec:v408-bibliography-citation-terminal-record-hygiene-cleanup}

\begin{quote}\small
\noindent\textbf{Status.}
This \QSCver{shifted.0.13} layer is a bibliography, citation, and terminal-record hygiene
cleanup. It does not change the shifted namespace shift, the backtranslation map,
the occurrence classes, display triples, tensor grammar, readout compatibility,
OP status, or the Historical Witness. It records how bibliographic and
terminal records should be read in the split bundle.
\end{quote}

\subsection{Bibliography role principle}
\label{subsec:v408-bibliography-role-principle}

\begin{definition}[Bibliography role principle]
\label{def:v408-bibliography-role-principle}
The bibliography and terminal self-citation records are metadata and scholarly
context. They do not define mathematical objects. Current mathematical
objects are governed by their definitions in the Active Core, especially
\[
 \Sigma_4,\qquad
 \mathrm B_4,\qquad
 \mathfrak D_4,\qquad
 \mathcal N_4.
\]
\end{definition}

\begin{remark}[Citation does not promote theorematic status]
\label{rem:v408-citation-does-not-promote-theorematic-status}
A citation, DOI record, or terminal self-citation can identify a release or
source, but it does not promote a historical statement to current theorematic
status. Promotion requires a current Active-Core normal-form record and status
guard.
\end{remark}

\subsection{Terminal record hygiene}
\label{subsec:v408-terminal-record-hygiene}

\begin{definition}[Terminal record hygiene rule]
\label{def:v408-terminal-record-hygiene-rule}
Terminal records should separate:
\begin{enumerate}[leftmargin=*]
 \item concept-level identifiers;
 \item version-level identifiers;
 \item prior public version identifiers;
 \item local cleanup-release-final state identifiers;
 \item bibliography entries and self-citations.
\end{enumerate}
No terminal record should be read as a new theorem unless it is explicitly
stated as a theorem in the Active Core.
\end{definition}

\begin{lemma}[Terminal records preserve metadata boundaries]
\label{lem:v408-terminal-records-preserve-metadata-boundaries}
The terminal record hygiene rule preserves the boundary between metadata and
mathematics.
\end{lemma}

\begin{proof}
The rule assigns identifiers and citations to metadata roles. It explicitly
prevents DOI records and self-citations from acting as theorem statements.
Therefore metadata remains separate from mathematical content.
\end{proof}

\subsection{Citation audit contract}
\label{subsec:v408-citation-audit-contract}

\begin{definition}[Citation audit contract]
\label{def:v408-citation-audit-contract}
A clean citation state for a local release-final cleanup version requires:
\begin{enumerate}[leftmargin=*]
 \item no undefined citations in the build log;
 \item no undefined references in the build log;
 \item the concept DOI and prior public version DOI to be readable as
 metadata;
 \item no unassigned DOI to be presented as a current version DOI;
 \item the Historical Witness to remain linked as proof memory rather than as
 a second active bibliography line.
\end{enumerate}
\end{definition}

\begin{remark}[Local release-final cleanup versions and DOI restraint]
\label{rem:v408-local-cleanup-release-final states-doi-restraint}
A local release-final cleanup version may mention the prior DOI and the concept DOI. It
should not claim a assigned version DOI until release-final packaging assigns one.
\end{remark}

\subsection{Bibliography and witness linkage}
\label{subsec:v408-bibliography-witness-linkage}

\begin{definition}[Bibliography--witness linkage]
\label{def:v408-bibliography-witness-linkage}
Historical citations and witness-preserved references may be used to reconstruct
the genealogy of an argument. They do not override the current shifted-namespace normal form.
If a historical citation supports a current statement, the current statement
must still be expressed in Active-Core language.
\end{definition}

\begin{proposition}[Bibliographic context is not a second governing layer]
\label{prop:v408-bibliographic-context-not-second-governing-layer}
Bibliographic context and Historical Witness context do not form a second
governing theorematic layer.
\end{proposition}

\begin{proof}
By Definition~\ref{def:v408-bibliography-witness-linkage}, historical citations
and witness references serve genealogy and reconstruction roles. Current
governance remains in the Active Core. Hence bibliographic context is not a
second governing layer.
\end{proof}

\subsection{Bibliography hygiene conservativity}
\label{subsec:v408-bibliography-hygiene-conservativity}

\begin{theorem}[Bibliography and citation hygiene conservativity]
\label{thm:v408-bibliography-citation-hygiene-conservativity}
The \QSCver{shifted.0.13} bibliography/citation/terminal-record hygiene cleanup is conservative
over \QSCver{shifted.0.7}.
\end{theorem}

\begin{proof}
The cleanup records how citations, DOI records, terminal records, and witness
references should be read. It does not modify the definitions of $\Sigma_4$,
$\mathrm B_4$, $\mathfrak D_4$, or $\mathcal N_4$. It does not alter tensor
grammar, readout compatibility, OP-status guards, or Historical Witness
content. Therefore it is conservative over \QSCver{shifted.0.7}.
\end{proof}

\begin{remark}[\QSCver{shifted.0.13} non-claim]
\label{rem:v408-non-claim}
Version~4.0.13 is a bibliography, citation, and terminal-record hygiene
release-final state. It does not add new physics, does not close OP~1/4 or OP~5(B),
does not activate primitive mixed tensors, and does not claim
Turing-completeness.
\end{remark}



\clearpage
\section{\QSCver{shifted.0.13} Glossary, Terminology, and Reader-Index Hygiene Cleanup}
\label{sec:v409-glossary-terminology-reader-index-hygiene-cleanup}

\begin{quote}\small
\noindent\textbf{Status.}
This \QSCver{shifted.0.13} layer is a glossary, terminology, and reader-index hygiene cleanup.
It does not change the shifted namespace shift, the backtranslation map, the
occurrence classes, display triples, tensor grammar, readout compatibility, OP
status, citations, metadata, or the Historical Witness. It supplies a compact
reader-facing terminology map for the current shifted line.
\end{quote}

\subsection{Terminology hygiene principle}
\label{subsec:v409-terminology-hygiene-principle}

\begin{definition}[Terminology hygiene principle]
\label{def:v409-terminology-hygiene-principle}
A reader-facing term is hygienic if it has exactly one active role in the
current shifted line, or if all of its roles are explicitly separated by context.
Terminology hygiene prevents a word such as ``tag'', ``mask'', ``witness'', or
``readout'' from silently changing mathematical role.
\end{definition}

\begin{remark}[Glossary entries are reader aids]
\label{rem:v409-glossary-entries-reader-aids}
The glossary entries in this layer are reader aids. They do not replace the
formal definitions in the Active Core and do not promote historical material to
current theorematic status.
\end{remark}

\subsection{Compact current-shifted glossary}
\label{subsec:v409-compact-current-shifted-glossary}

\begin{center}
\scriptsize
\setlength{\tabcolsep}{2pt}
\begin{tabular}{@{}p{0.18\textwidth}P{0.35\textwidth}P{0.34\textwidth}@{}}
\toprule
Term & Current shifted role & Guard against confusion\\
\midrule
Active Core & governing current line & not the Historical Witness\\
Historical Witness & proof memory and genealogy & not a second active theory line\\
QAM & language/router/readout set layer & not a new physical law by itself\\
MML & structural/router-machine layer & not a replacement for proof status\\
MMA & matrix/readout machinery layer & not an OP-closure certificate\\
HC & Heisenberg compensator/boundary normalizer & not redefined by the global shifted shift\\
tag & inhabited active QAM/router label & not a packed mask integer\\
route payload & ordered finite tag/readout route & not an unordered support set\\
support set & unordered finite support entries & shifted before mask recomputation\\
packed mask & integer encoding of support & not shifted by adding one\\
tensor-index payload & QAM/router payload in an index position & not the tensor component itself\\
readout target & physical/mathematical target of a package & not changed by address migration\\
guard & non-claim/status protection record & not optional prose decoration\\
OP status & open/closed theorematic state & not changed by namespace migration\\
\bottomrule
\end{tabular}
\end{center}

\subsection{Acronym and abbreviation map}
\label{subsec:v409-acronym-abbreviation-map}

\begin{definition}[Acronym map]
\label{def:v409-acronym-map}
The following acronyms are read in the current Active Core as follows:
\[
\begin{array}{ll}
\mathrm{QAM} & \text{Quantum Analysis Mengenlehre / language and set-readout layer},\\
\mathrm{MML} & \text{Millennium Matrix / structural router-machine layer},\\
\mathrm{MMA} & \text{Heisenberg Matrix / matrix-readout machinery layer},\\
\mathrm{HC} & \text{Heisenberg compensator / boundary normalizer},\\
\mathrm{OP} & \text{open problem or OP-facing package, depending on context}.
\end{array}
\]
\end{definition}

\begin{remark}[Acronyms do not carry closure]
\label{rem:v409-acronyms-do-not-carry-closure}
An acronym can identify a layer or tool family, but it does not by itself close
an OP or certify a physical readout.
\end{remark}

\subsection{Reader-index roles}
\label{subsec:v409-reader-index-roles}

\begin{definition}[Reader-index role]
\label{def:v409-reader-index-role}
A reader-index role is a stable navigation label attached to a term. The
current cleanup line uses four reader-index roles:
\[
 \text{governing},\qquad
 \text{routing},\qquad
 \text{witness},\qquad
 \text{guard}.
\]
Governing terms state the active line; routing terms explain transport or
display; witness terms preserve genealogy; guard terms protect non-claims and
OP status.
\end{definition}

\begin{lemma}[Reader-index roles prevent role drift]
\label{lem:v409-reader-index-roles-prevent-role-drift}
The reader-index roles prevent glossary terms from silently switching between
active, routing, witness, and guard functions.
\end{lemma}

\begin{proof}
Each role has a specified reading function. If a term is used as a governing
term, it must belong to the active line. If it is used as witness material, it
does not govern. If it is used as a guard, it protects status. If it is used
as routing, it explains display or transport. Thus the role determines the
reading and prevents silent drift.
\end{proof}

\subsection{Terminology conservativity}
\label{subsec:v409-terminology-conservativity}

\begin{theorem}[Glossary and terminology hygiene conservativity]
\label{thm:v409-glossary-terminology-hygiene-conservativity}
The \QSCver{shifted.0.13} glossary/terminology/reader-index hygiene cleanup is conservative
over \QSCver{shifted.0.8}.
\end{theorem}

\begin{proof}
The cleanup adds reader-facing terminology maps and index roles. It does not
modify the definitions of $\Sigma_4$, $\mathrm B_4$, $\mathfrak D_4$, or
$\mathcal N_4$. It does not alter tensor grammar, readout compatibility,
citation records, metadata, Historical Witness content, or OP-status guards.
Therefore it is conservative over \QSCver{shifted.0.8}.
\end{proof}

\begin{remark}[\QSCver{shifted.0.13} non-claim]
\label{rem:v409-non-claim}
Version~4.0.13 is a glossary, terminology, and reader-index hygiene release-final state.
It does not add new physics, does not close OP~1/4 or OP~5(B), does not
activate primitive mixed tensors, and does not claim Turing-completeness.
\end{remark}



\clearpage
\section{\QSCver{shifted.0.13} Cleanup Consolidation and Local Release-Final Preparation}
\label{sec:v4010-cleanup-consolidation-local-release-release-final state-preparation}

\begin{quote}\small
\noindent\textbf{Status.}
This \QSCver{shifted.0.13} layer consolidates the local cleanup sequence into one
release-release-final state-preparation layer. It does not add new mathematics, does not
close OP~1/4 or OP~5(B), does not activate primitive mixed tensors, and does
not change the \QSCver{shifted.0.0} global QAM Hilbert-shift. It only records that the
cleanup sequence has become locally stable enough for a possible small
DOI-bearing cleanup release.
\end{quote}

\subsection{Cleanup consolidation principle}
\label{subsec:v4010-cleanup-consolidation-principle}

\begin{definition}[Cleanup consolidation principle]
\label{def:v4010-cleanup-consolidation-principle}
A cleanup consolidation layer is a reader-facing stabilization layer which
collects prior cleanup work without changing the mathematical content. For the
current line, it consolidates:
\[
\begin{array}{ll}
\text{\QSCver{shifted.0.1}} & \text{historical remark triage},\\
\text{\QSCver{shifted.0.2}} & \text{reader navigation},\\
\text{\QSCver{shifted.0.3}} & \text{notation and symbol hygiene},\\
\text{\QSCver{shifted.0.4}} & \text{table/display/example normalization},\\
\text{\QSCver{shifted.0.5}} & \text{Historical Witness linkage},\\
\text{\QSCver{shifted.0.6}} & \text{clean reader bundle},\\
\text{\QSCver{shifted.0.7}} & \text{frontmatter and DOI-record hygiene},\\
\text{\QSCver{shifted.0.8}} & \text{bibliography/citation/terminal-record hygiene},\\
\text{\QSCver{shifted.0.9}} & \text{glossary/terminology/reader-index hygiene}.
\end{array}
\]
\end{definition}

\begin{remark}[Consolidation is not compression by deletion]
\label{rem:v4010-consolidation-not-deletion}
The consolidation does not delete proof memory. It leaves the Historical
Witness preserved and treats cleanup as reader-facing stabilization rather than
content removal.
\end{remark}

\subsection{Local release-release-final state readiness checklist}
\label{subsec:v4010-local-release-release-final state-readiness-checklist}

\begin{definition}[Local cleanup release-release-final state readiness]
\label{def:v4010-local-cleanup-release-release-final state-readiness}
The cleanup-line is locally ready for release-release-final state packaging when
the following conditions hold:
\begin{enumerate}[leftmargin=*]
 \item the Active Core compiles without undefined references or citations;
 \item the Historical Witness compiles without undefined references or
 citations;
 \item no duplicate labels are reported;
 \item overfull hbox/vbox warnings are absent in both PDFs;
 \item the title and concept DOI remain stable;
 \item the prior public version DOI is recorded as prior, not as a new DOI;
 \item all non-claim guards remain visible.
\end{enumerate}
\end{definition}

\begin{lemma}[Readiness checklist is metadata-level]
\label{lem:v4010-readiness-checklist-metadata-level}
The local release-release-final state readiness checklist is a metadata and build-audit
condition, not a theorematic closure condition.
\end{lemma}

\begin{proof}
The checklist speaks about build state, title/DOI metadata, duplicate labels,
warnings, and non-claim guards. It does not add proof certificates or change
the status record of any mathematical object. Therefore it is metadata-level.
\end{proof}

\subsection{Current invariant freeze}
\label{subsec:v4010-current-invariant-freeze}

\begin{definition}[Current cleanup invariant freeze]
\label{def:v4010-current-cleanup-invariant-freeze}
Throughout the cleanup sequence, the following invariants remain fixed:
\[
\begin{array}{ll}
\Sigma_4(k)=k+1, & \mathrm B_4(a)=a-1\quad(a\geq1),\\[0.2em]
0_{\mathrm{shift}}=\varnothing_{\mathrm{active}}, & 0_{\mathrm{mask}}\to1_{\mathrm{mask}},\\[0.2em]
\mathfrak D_4(X), & \mathcal N_4(X),\\[0.2em]
\text{OP~1/4 open}, & \text{OP~5(B) open}.
\end{array}
\]
\end{definition}

\begin{proposition}[Cleanup consolidation preserves the invariant freeze]
\label{prop:v4010-cleanup-consolidation-preserves-invariant-freeze}
The \QSCver{shifted.0.13} cleanup consolidation preserves the invariant freeze of
Definition~\ref{def:v4010-current-cleanup-invariant-freeze}.
\end{proposition}

\begin{proof}
The consolidation layer only records the cleanup history and readiness
conditions. It does not modify the shift, backtranslation, reserved zero,
HC-boundary shift, display triple, normal form, or OP-status guards. Hence all
listed invariants are preserved.
\end{proof}

\subsection{Release packaging boundary}
\label{subsec:v4010-release-packaging-boundary}

\begin{definition}[Release packaging boundary]
\label{def:v4010-release-packaging-boundary}
Release packaging begins only when a assigned Zenodo version DOI is explicitly
assigned. Before that assignment, \QSCver{shifted.0.13} remains a local cleanup consolidation
release-final state with the prior public version DOI recorded only as prior.
\end{definition}

\begin{remark}[What a final cleanup release would do]
\label{rem:v4010-final-cleanup-release-would-do}
A final cleanup release would not add new theory. It would insert the
assigned version DOI, normalize release-final file names, regenerate audit reports,
checksums, and the Zenodo description, and preserve the Active Core/Historical
Witness split.
\end{remark}

\subsection{Cleanup consolidation conservativity}
\label{subsec:v4010-cleanup-consolidation-conservativity}

\begin{theorem}[Cleanup consolidation conservativity]
\label{thm:v4010-cleanup-consolidation-conservativity}
The \QSCver{shifted.0.13} cleanup consolidation is conservative over \QSCver{shifted.0.0} release-final and
over the local \QSCver{shifted.0.1}--\QSCver{shifted.0.9} cleanup sequence.
\end{theorem}

\begin{proof}
The consolidation summarizes the cleanup line and records release-release-final state
readiness criteria. It does not alter any current shifted-namespace mathematical definition,
tensor grammar, readout compatibility check, citation role, metadata role,
Historical Witness content, or OP-status guard. Therefore it is conservative
over \QSCver{shifted.0.0} and over \QSCver{shifted.0.1}--\QSCver{shifted.0.9}.
\end{proof}

\begin{remark}[\QSCver{shifted.0.13} non-claim]
\label{rem:v4010-non-claim}
Version~4.0.13 is a cleanup consolidation release-final state. It does not add new
physics, does not close OP~1/4 or OP~5(B), does not activate primitive mixed
tensors, and does not claim Turing-completeness.
\end{remark}



\clearpage
\section{\QSCver{shifted.0.13} Label, Reference, and Counter Hygiene Cleanup}
\label{sec:v4011-label-reference-counter-hygiene-cleanup}

\begin{quote}\small
\noindent\textbf{Status.}
This \QSCver{shifted.0.13} layer is a label, reference, and counter hygiene cleanup. It does
not change the shifted namespace shift, the backtranslation map, the occurrence
classes, display triples, tensor grammar, readout compatibility, OP status,
citations, metadata, or the Historical Witness. It records the audit rule for
labels, cross-references, theorem counters, and reader-visible navigation.
\end{quote}

\subsection{Label hygiene principle}
\label{subsec:v4011-label-hygiene-principle}

\begin{definition}[Label hygiene principle]
\label{def:v4011-label-hygiene-principle}
A label is hygienic if it is unique, stable under cleanup edits, and
semantically aligned with the object it names. In the cleanup-line,
labels should identify the local role of an object, for example section,
definition, lemma, theorem, proposition, remark, table, or audit block.
\end{definition}

\begin{remark}[Labels are addresses, not mathematical content]
\label{rem:v4011-labels-are-addresses}
A \LaTeX{} label is a document address. It does not define the mathematical
object itself. Mathematical content remains in the corresponding definition,
lemma, theorem, proposition, or normal-form record.
\end{remark}

\subsection{Reference hygiene contract}
\label{subsec:v4011-reference-hygiene-contract}

\begin{definition}[Reference hygiene contract]
\label{def:v4011-reference-hygiene-contract}
A clean local release-final cleanup version should satisfy the following reference contract:
\begin{enumerate}[leftmargin=*]
 \item no undefined references in the build log;
 \item no undefined citations in the build log;
 \item no multiply defined labels in the build log;
 \item no theorem statement should depend on an unresolved reference;
 \item historical references should not override current Active-Core status.
\end{enumerate}
\end{definition}

\begin{lemma}[Reference hygiene is build-level, not theorematic]
\label{lem:v4011-reference-hygiene-build-level}
Reference hygiene is a build-level and reader-navigation condition, not a new
theorematic closure condition.
\end{lemma}

\begin{proof}
The reference contract concerns build-log state and reader addressability. It
does not change any mathematical definition, readout target, guard, or OP-status
record. Therefore it is build-level rather than theorematic.
\end{proof}

\subsection{Counter and theorem-environment hygiene}
\label{subsec:v4011-counter-theorem-environment-hygiene}

\begin{definition}[Counter hygiene rule]
\label{def:v4011-counter-hygiene-rule}
Theorem, lemma, proposition, definition, and remark counters are reader-facing
navigation devices. Their role is to keep the document addressable. A cleanup
pass should not use counter renumbering to imply a change of theorematic status.
\end{definition}

\begin{proposition}[Counter changes are not status changes]
\label{prop:v4011-counter-changes-not-status-changes}
A reader-visible counter or page-number change caused by cleanup is not a change
of theorematic status.
\end{proposition}

\begin{proof}
The theorematic status is recorded by the mathematical statement and its guard
or status record, not by its page number or counter position. Cleanup may add
reader-facing sections and thereby shift page numbers. Such shifts do not
alter status records. Therefore counter or page-number changes are not status
changes.
\end{proof}

\subsection{Cross-reference role separation}
\label{subsec:v4011-cross-reference-role-separation}

\begin{definition}[Cross-reference roles]
\label{def:v4011-cross-reference-roles}
Cross-references are read in one of four roles:
\[
 \text{governing reference},\qquad
 \text{routing reference},\qquad
 \text{witness reference},\qquad
 \text{audit reference}.
\]
A governing reference points to current Active-Core content. A routing
reference explains transport, display, or navigation. A witness reference
points to proof memory or genealogy. An audit reference points to build,
metadata, or hygiene state.
\end{definition}

\begin{remark}[Witness references remain non-governing]
\label{rem:v4011-witness-references-non-governing}
A reference to the Historical Witness remains a witness reference unless the
Active Core explicitly reintroduces the content through a current shifted display
triple, normal form, or status record.
\end{remark}

\subsection{Label/reference hygiene conservativity}
\label{subsec:v4011-label-reference-hygiene-conservativity}

\begin{theorem}[Label/reference/counter hygiene conservativity]
\label{thm:v4011-label-reference-counter-hygiene-conservativity}
The \QSCver{shifted.0.13} label/reference/counter hygiene cleanup is conservative over
\QSCver{shifted.0.10}.
\end{theorem}

\begin{proof}
The cleanup records build-addressing and navigation rules. It does not modify
the definitions of $\Sigma_4$, $\mathrm B_4$, $\mathfrak D_4$, or
$\mathcal N_4$. It does not alter tensor grammar, readout compatibility,
citation records, metadata, Historical Witness content, or OP-status guards.
Therefore it is conservative over \QSCver{shifted.0.10}.
\end{proof}

\begin{remark}[\QSCver{shifted.0.13} non-claim]
\label{rem:v4011-non-claim}
Version~4.0.13 is a label, reference, and counter hygiene release-final state. It does
not add new physics, does not close OP~1/4 or OP~5(B), does not activate
primitive mixed tensors, and does not claim Turing-completeness.
\end{remark}



\clearpage
\section{\QSCver{shifted.0.13} Layout, TOC, and Render Hygiene Cleanup}
\label{sec:v4012-layout-toc-render-hygiene-cleanup}

\begin{quote}\small
\noindent\textbf{Status.}
This \QSCver{shifted.0.13} layer is a layout, table-of-contents, and render hygiene cleanup.
It does not change the shifted namespace shift, the backtranslation map, the
occurrence classes, display triples, tensor grammar, readout compatibility, OP
status, citations, metadata, labels, counters, or the Historical Witness. It
records the visual audit rule for the split bundle.
\end{quote}

\subsection{Layout hygiene principle}
\label{subsec:v4012-layout-hygiene-principle}

\begin{definition}[Layout hygiene principle]
\label{def:v4012-layout-hygiene-principle}
A local release-final cleanup version is layout-hygienic if its generated PDFs are readable
as rendered documents: no text is visibly clipped, no major text blocks overlap,
table headings remain legible, and reader-navigation pages are visually usable.
\end{definition}

\begin{remark}[Layout is not theorematic content]
\label{rem:v4012-layout-not-theorematic-content}
Layout hygiene improves readability and auditability. It does not change
definitions, lemmas, theorems, OP status, or physical readout status.
\end{remark}

\subsection{TOC readability contract}
\label{subsec:v4012-toc-readability-contract}

\begin{definition}[TOC readability contract]
\label{def:v4012-toc-readability-contract}
The table of contents is reader-hygienic if section and subsection entries are
readable, numbering boxes do not collide with titles, and cleanup-version
headings remain navigable. Long cleanup headings may be verbose in the source
but must remain readable in the rendered PDF.
\end{definition}

\begin{lemma}[TOC readability is a navigation condition]
\label{lem:v4012-toc-readability-navigation-condition}
The TOC readability contract is a navigation condition, not a mathematical
condition.
\end{lemma}

\begin{proof}
The table of contents addresses document locations. It does not define or
modify mathematical objects. Therefore improving TOC readability changes only
navigation.
\end{proof}

\subsection{Render audit contract}
\label{subsec:v4012-render-audit-contract}

\begin{definition}[Render audit contract]
\label{def:v4012-render-audit-contract}
A local render audit should include at least the title page, early navigation
pages, at least one representative middle page, and terminal pages of both the
Active Core and Historical Witness. The purpose is to detect visible clipping,
overlap, unreadable page headings, and severe table breakage.
\end{definition}

\begin{remark}[Contact sheet role]
\label{rem:v4012-contact-sheet-role}
The render contact sheet is an audit artifact. It is not part of the
mathematical proof. It gives a compact visual sample of the generated PDFs.
\end{remark}

\subsection{Warning hygiene}
\label{subsec:v4012-warning-hygiene}

\begin{definition}[Warning hygiene contract]
\label{def:v4012-warning-hygiene-contract}
For the current cleanup line, the preferred local build state is:
\[
\begin{array}{ll}
\text{undefined references} & =0,\\
\text{undefined citations} & =0,\\
\text{multiply-defined labels} & =0,\\
\text{overfull hbox warnings} & =0,\\
\text{overfull vbox warnings} & =0.
\end{array}
\]
This is a local document-hygiene target rather than a theorematic claim.
\end{definition}

\begin{proposition}[Warning hygiene is conservative]
\label{prop:v4012-warning-hygiene-conservative}
Achieving the preferred warning hygiene state does not alter mathematical
content.
\end{proposition}

\begin{proof}
Warnings and references concern build quality, cross-reference resolution, and
rendering. They do not change the definitions of $\Sigma_4$, $\mathrm B_4$,
$\mathfrak D_4$, or $\mathcal N_4$, and they do not modify OP-status guards.
Therefore warning hygiene is conservative.
\end{proof}

\subsection{Layout/render hygiene conservativity}
\label{subsec:v4012-layout-render-hygiene-conservativity}

\begin{theorem}[Layout/TOC/render hygiene conservativity]
\label{thm:v4012-layout-toc-render-hygiene-conservativity}
The \QSCver{shifted.0.13} layout/TOC/render hygiene cleanup is conservative over \QSCver{shifted.0.11}.
\end{theorem}

\begin{proof}
The cleanup records visual audit, TOC readability, and warning-hygiene rules.
It does not modify any current shifted-namespace mathematical definition, tensor grammar,
readout compatibility check, citation role, metadata role, label role,
Historical Witness content, or OP-status guard. Therefore it is conservative
over \QSCver{shifted.0.11}.
\end{proof}

\begin{remark}[\QSCver{shifted.0.13} non-claim]
\label{rem:v4012-non-claim}
Version~4.0.13 is a layout, TOC, and render hygiene release-final state. It does not add
new physics, does not close OP~1/4 or OP~5(B), does not activate primitive
mixed tensors, and does not claim Turing-completeness.
\end{remark}



\clearpage
\section{\QSCver{shifted.0.13} Source, Build, and Artifact Reproducibility Hygiene Cleanup}
\label{sec:v4013-source-build-artifact-reproducibility-hygiene-cleanup}

\begin{quote}\small
\noindent\textbf{Status.}
This \QSCver{shifted.0.13} layer is a source, build, and artifact reproducibility hygiene
cleanup. It does not change the shifted namespace shift, the backtranslation map,
the occurrence classes, display triples, tensor grammar, readout compatibility,
OP status, citations, metadata, labels, counters, layout contract, or the
Historical Witness. It records how source files, generated PDFs, manifests,
checksums, and bundles should be interpreted.
\end{quote}

\subsection{Source reproducibility principle}
\label{subsec:v4013-source-reproducibility-principle}

\begin{definition}[Source reproducibility principle]
\label{def:v4013-source-reproducibility-principle}
A local release-final cleanup version is source-reproducible if the supplied TEX files are
sufficient to rebuild the Active Core and Historical Witness PDFs in the same
logical document state, modulo ordinary build-environment differences such as
PDF timestamps or engine-specific metadata.
\end{definition}

\begin{remark}[Source is governing over generated artifacts]
\label{rem:v4013-source-governing-over-generated-artifacts}
For local cleanup work, the TEX source is the editable governing artifact. The
PDFs are rendered outputs. If a rendered artifact and source artifact appear to
disagree, the source and build log must be inspected before any mathematical
interpretation is changed.
\end{remark}

\subsection{Artifact classes}
\label{subsec:v4013-artifact-classes}

\begin{definition}[Artifact classes]
\label{def:v4013-artifact-classes}
The local cleanup workflow separates artifacts into the following classes:
\begin{enumerate}[leftmargin=*]
 \item \textbf{source artifacts:} TEX files and source-side scan reports;
 \item \textbf{rendered artifacts:} Active Core and Historical Witness PDFs;
 \item \textbf{audit artifacts:} build logs, audit reports, render contact
 sheets, and scan reports;
 \item \textbf{packaging artifacts:} manifest files, checksum files, full
 bundles, and TEX-only ZIPs.
\end{enumerate}
\end{definition}

\begin{lemma}[Artifact classes prevent role drift]
\label{lem:v4013-artifact-classes-prevent-role-drift}
The artifact classes prevent rendered, audit, and packaging artifacts from
being mistaken for new mathematical content.
\end{lemma}

\begin{proof}
Each artifact class has an assigned role. Source artifacts contain editable
content; rendered artifacts show compiled output; audit artifacts document
quality checks; packaging artifacts support transfer and verification. None of
the non-source classes is a new theorematic layer. Therefore role drift is
prevented.
\end{proof}

\subsection{Manifest and checksum hygiene}
\label{subsec:v4013-manifest-checksum-hygiene}

\begin{definition}[Manifest and checksum hygiene rule]
\label{def:v4013-manifest-checksum-hygiene-rule}
A local cleanup bundle should contain a manifest and checksum file. The
manifest records which files belong to the package. The checksum file records
file integrity. Neither manifest nor checksum file changes mathematical
content.
\end{definition}

\begin{proposition}[Checksums certify files, not theorems]
\label{prop:v4013-checksums-certify-files-not-theorems}
A checksum certifies the byte identity of a file, not the truth of a theorem.
\end{proposition}

\begin{proof}
A checksum is a hash of file bytes. It can detect whether a file has changed,
but it does not inspect or prove mathematical statements. Therefore it
certifies file integrity only.
\end{proof}

\subsection{Build provenance contract}
\label{subsec:v4013-build-provenance-contract}

\begin{definition}[Build provenance contract]
\label{def:v4013-build-provenance-contract}
A clean local build record should state:
\begin{enumerate}[leftmargin=*]
 \item source TEX file names;
 \item output PDF file names;
 \item page counts;
 \item undefined reference and citation counts;
 \item duplicate-label counts;
 \item overfull hbox and vbox counts;
 \item fatal or emergency error counts;
 \item whether a render contact sheet was produced.
\end{enumerate}
\end{definition}

\begin{remark}[Build provenance is audit support]
\label{rem:v4013-build-provenance-audit-support}
Build provenance supports reproducibility and audit. It does not replace the
mathematical proof and does not close any OP.
\end{remark}

\subsection{TEX-only bundle contract}
\label{subsec:v4013-tex-only-bundle-contract}

\begin{definition}[TEX-only bundle contract]
\label{def:v4013-tex-only-bundle-contract}
A TEX-only bundle is intended for manual integration and local rebuilding. It
should contain the Active Core source, the Historical Witness source, and the
minimal supporting hygiene records needed to identify the release-final state. It should
not be treated as a DOI-bearing release unless release-final packaging assigns
a assigned version DOI.
\end{definition}

\begin{lemma}[TEX-only bundle preserves manual-integration workflow]
\label{lem:v4013-tex-only-bundle-preserves-manual-integration-workflow}
The TEX-only bundle contract preserves the user's manual-integration workflow.
\end{lemma}

\begin{proof}
The contract supplies directly editable TEX files and minimal identifiers. It
does not require repository-side mutation, automatic patch application, or a
new DOI. Hence it preserves manual integration.
\end{proof}

\subsection{Reproducibility conservativity}
\label{subsec:v4013-reproducibility-conservativity}

\begin{theorem}[Source/build/artifact hygiene conservativity]
\label{thm:v4013-source-build-artifact-hygiene-conservativity}
The \QSCver{shifted.0.13} source/build/artifact reproducibility cleanup is conservative over
\QSCver{shifted.0.12}.
\end{theorem}

\begin{proof}
The cleanup records artifact roles, manifest/checksum rules, build provenance,
and TEX-only bundle expectations. It does not modify the definitions of
$\Sigma_4$, $\mathrm B_4$, $\mathfrak D_4$, or $\mathcal N_4$. It does not
alter tensor grammar, readout compatibility, citation records, metadata roles,
labels, counters, layout rules, Historical Witness content, or OP-status
guards. Therefore it is conservative over \QSCver{shifted.0.12}.
\end{proof}

\begin{remark}[\QSCver{shifted.0.13} non-claim]
\label{rem:v4013-non-claim}
Version~4.0.13 is a source, build, and artifact reproducibility hygiene
release-final state. It does not add new physics, does not close OP~1/4 or OP~5(B),
does not activate primitive mixed tensors, and does not claim
Turing-completeness.
\end{remark}



\clearpage
\section{\QSCver{shifted.0.14} Local Release-Readiness, Upload, and Final-Packaging Hygiene Cleanup}
\label{sec:v4014-local-release-readiness-upload-final-packaging-hygiene-cleanup}

\begin{quote}\small
\noindent\textbf{Status.}
This \QSCver{shifted.0.14} layer is a local release-readiness and upload-hygiene cleanup. It
records the assigned version DOI 10.5281/zenodo.19811662, does not change the shifted namespace shift, does not
change the Historical Witness, and does not change any OP status. It records
which artifacts would be used for a small DOI-bearing cleanup release if a new
Zenodo version DOI is later assigned.
\end{quote}

\subsection{Local release-readiness principle}
\label{subsec:v4014-local-release-readiness-principle}

\begin{definition}[Local release-readiness principle]
\label{def:v4014-local-release-readiness-principle}
A local release-final cleanup version is release-ready only in the limited sense that its
source, PDF, audit, manifest, checksum, and bundle artifacts are coherent and
can be promoted to release-final packaging after the version DOI 10.5281/zenodo.19811662 is assigned.
Release-readiness is not a theorematic status and is not an OP-closure
certificate.
\end{definition}

\begin{remark}[release-final assigned in this layer]
\label{rem:v4014-version-doi-assigned-release-final}
This release-final layer records the assigned Zenodo version DOI 10.5281/zenodo.19811662. The prior public
version DOI remains the \QSCver{shifted.0.0} release DOI until explicit release-final
packaging assigns a assigned version DOI.
\end{remark}

\subsection{Upload artifact selection}
\label{subsec:v4014-upload-artifact-selection}

\begin{definition}[Recommended upload artifact set]
\label{def:v4014-recommended-upload-artifact-set}
For a small cleanup release, the recommended upload artifact set is:
\begin{enumerate}[leftmargin=*]
 \item Active Core PDF;
 \item Historical Witness PDF;
 \item full release bundle ZIP;
 \item TEX-only ZIP, if source-side manual integration is desired;
 \item release notes or Zenodo description text;
 \item checksum file or checksum-bearing manifest.
\end{enumerate}
\end{definition}

\begin{remark}[Candidate artifacts versus release-final artifacts]
\label{rem:v4014-release-final state-artifacts-versus-release-final-artifacts}
Candidate artifacts should not be uploaded as release-final artifacts unless
their file names, metadata, DOI records, manifests, checksums, and Zenodo
description have been normalized for release-final.
\end{remark}

\subsection{Zenodo description draft contract}
\label{subsec:v4014-zenodo-description-draft-contract}

\begin{definition}[Zenodo description draft contract]
\label{def:v4014-zenodo-description-draft-contract}
A Zenodo description draft for a local release-final cleanup version should state:
\begin{enumerate}[leftmargin=*]
 \item the document title;
 \item the cleanup version;
 \item that the update is reader/document hygiene, not new theory;
 \item the concept DOI;
 \item the prior public version DOI;
 \item that the assigned version DOI is assigned in this release-final packaging as 10.5281/zenodo.19811662;
 \item the non-claims for OP~1/4, OP~5(B), primitive mixed tensors, and
 Turing-completeness.
\end{enumerate}
\end{definition}

\begin{lemma}[Description draft is metadata]
\label{lem:v4014-description-draft-is-metadata}
A Zenodo description draft is metadata and does not change mathematical content.
\end{lemma}

\begin{proof}
The description draft summarizes the package, version state, DOI state, and
non-claims. It does not alter definitions, proofs, normal forms, readout
guards, or OP-status records. Hence it is metadata only.
\end{proof}

\subsection{Final-packaging boundary}
\label{subsec:v4014-final-packaging-boundary}

\begin{definition}[Final-packaging boundary]
\label{def:v4014-final-packaging-boundary}
Final packaging begins only after a assigned Zenodo version DOI is provided. The
final-packaging pass may then:
\begin{enumerate}[leftmargin=*]
 \item replace release-final state wording by release-final wording where appropriate;
 \item insert the assigned version DOI in metadata, manifest, notes, and Zenodo
 description;
 \item normalize file names to release-final names;
 \item rebuild the Active Core and Historical Witness PDFs;
 \item regenerate audit reports, checksums, and bundles.
\end{enumerate}
\end{definition}

\begin{proposition}[Final packaging is conservative when DOI-only]
\label{prop:v4014-final-packaging-conservative-when-doi-only}
If final packaging only inserts DOI and release-final metadata, rebuilds
artifacts, and regenerates audits, then it is conservative over the local
release-final cleanup version.
\end{proposition}

\begin{proof}
DOI insertion and release-final artifact regeneration are metadata and build
operations. They do not change $\Sigma_4$, $\mathrm B_4$, $\mathfrak D_4$,
$\mathcal N_4$, tensor grammar, readout guards, Historical Witness content, or
OP-status records. Therefore such packaging is conservative.
\end{proof}

\subsection{Release-readiness conservativity}
\label{subsec:v4014-release-readiness-conservativity}

\begin{theorem}[Release-readiness and upload-hygiene conservativity]
\label{thm:v4014-release-readiness-upload-hygiene-conservativity}
The \QSCver{shifted.0.14} release-readiness/upload-hygiene cleanup is conservative over
\QSCver{shifted.0.13}.
\end{theorem}

\begin{proof}
The cleanup records upload artifact selection, Zenodo description drafting, and
the final-packaging boundary. It does not modify current shifted mathematical
definitions, tensor grammar, readout compatibility checks, artifact roles,
build rules, Historical Witness content, or OP-status guards. Therefore it is
conservative over \QSCver{shifted.0.13}.
\end{proof}

\begin{remark}[\QSCver{shifted.0.14} non-claim]
\label{rem:v4014-non-claim}
Version~4.0.14 is a local release-readiness and upload-hygiene release-final state. It
does not add new physics, does not close OP~1/4 or OP~5(B), does not activate
primitive mixed tensors, does not claim Turing-completeness, and records the assigned Zenodo version DOI 10.5281/zenodo.19811662.
\end{remark}



\clearpage
\section{\QSCver{shifted.0.15} Local Release-Final Freeze and DOI-Request Packet Hygiene}
\label{sec:v4015-local-release-release-final state-freeze-doi-request-packet-hygiene}

\begin{quote}\small
\noindent\textbf{Status.}
This \QSCver{shifted.0.15} layer freezes the local cleanup line as a release-release-final state
preparation state. It records the assigned version DOI 10.5281/zenodo.19811662, does not alter the shifted
namespace shift, does not modify the Historical Witness, and does not change
any OP status. Its only purpose is to mark the local cleanup sequence as ready
for a possible DOI-request and release-final packaging pass.
\end{quote}

\subsection{Local release-release-final state freeze principle}
\label{subsec:v4015-local-release-release-final state-freeze-principle}

\begin{definition}[Local release-release-final state freeze principle]
\label{def:v4015-local-release-release-final state-freeze-principle}
A release-final state is a document-control state in which the
source, rendered PDFs, audit reports, manifests, checksums, and upload
description are internally coherent and should not receive new mathematical
content before a release-final DOI pass.
\end{definition}

\begin{remark}[Freeze is not mathematical closure]
\label{rem:v4015-freeze-not-mathematical-closure}
A release-release-final state freeze is a workflow boundary. It is not a proof of OP
closure, not a physical promotion criterion, and not a replacement for any
status record inside the shifted-namespace normal form.
\end{remark}

\subsection{DOI-request packet}
\label{subsec:v4015-doi-request-packet}

\begin{definition}[DOI-request packet]
\label{def:v4015-doi-request-packet}
A DOI-request packet for a small cleanup release should contain:
\begin{enumerate}[leftmargin=*]
 \item the frozen Active Core PDF;
 \item the frozen Historical Witness PDF;
 \item the full bundle ZIP;
 \item the TEX-only ZIP;
 \item the release-final Zenodo description;
 \item the audit report;
 \item the manifest and checksum file.
\end{enumerate}
The packet is still local until a assigned Zenodo version DOI is assigned.
\end{definition}

\subsection{Allowed release-final pass}
\label{subsec:v4015-allowed-release-final-pass}

\begin{definition}[Allowed release-final pass]
\label{def:v4015-allowed-release-final-pass}
After a assigned Zenodo version DOI is supplied, the allowed release-final pass may:
\begin{enumerate}[leftmargin=*]
 \item insert the version DOI in frontmatter, notes, manifest, checksums, audit
 metadata, and Zenodo description;
 \item replace local release-final state/freeze wording by release-final wording where
 appropriate;
 \item rename files from local/freeze/release-final state names to release-final names;
 \item rebuild the Active Core and Historical Witness PDFs;
 \item regenerate the full bundle, TEX-only ZIP, checksums, audit report, and
 final Zenodo description.
\end{enumerate}
It may not add new theorems, change OP status, or alter shifted normal-form content.
\end{definition}

\begin{proposition}[DOI pass is conservative under the allowed boundary]
\label{prop:v4015-doi-pass-conservative-under-allowed-boundary}
A release-final pass satisfying Definition~\ref{def:v4015-allowed-release-final-pass}
is conservative over the local \QSCver{shifted.0.15} freeze.
\end{proposition}

\begin{proof}
The allowed pass is restricted to DOI insertion, wording normalization, file
renaming, rebuilding, and audit regeneration. It explicitly excludes new
theorems, OP-status changes, and shifted normal-form changes. Therefore it is
conservative over the frozen local release-final state.
\end{proof}

\subsection{Frozen non-claim boundary}
\label{subsec:v4015-frozen-non-claim-boundary}

\begin{definition}[Frozen non-claim boundary]
\label{def:v4015-frozen-non-claim-boundary}
The release-final state preserves the following non-claims:
\[
\begin{array}{ll}
\text{OP~1/4} & \text{remains open},\\
\text{OP~5(B)} & \text{remains open},\\
\text{primitive mixed tensors} & \text{not activated},\\
\text{Turing-completeness} & \text{not claimed},\\
\text{standard physical calculations} & \text{not replaced}.
\end{array}
\]
\end{definition}

\subsection{Freeze conservativity}
\label{subsec:v4015-freeze-conservativity}

\begin{theorem}[Local release-release-final state freeze conservativity]
\label{thm:v4015-local-release-release-final state-freeze-conservativity}
The \QSCver{shifted.0.15} release-final state is conservative over the \QSCver{shifted.0.1}--\QSCver{shifted.0.14}
cleanup sequence and over the public \QSCver{shifted.0.0} release-final mathematical content.
\end{theorem}

\begin{proof}
The freeze adds document-control and DOI-request packet rules. It does not
modify $\Sigma_4$, $\mathrm B_4$, $\mathfrak D_4$, $\mathcal N_4$, tensor
grammar, readout compatibility, artifact roles, build rules, Historical Witness
content, or OP-status guards. Therefore it is conservative over the cleanup
sequence and over the public \QSCver{shifted.0.0} mathematical content.
\end{proof}

\begin{remark}[\QSCver{shifted.0.15} non-claim]
\label{rem:v4015-non-claim}
Version~4.0.15 is a release-final state. It does not add new
physics, does not close OP~1/4 or OP~5(B), does not activate primitive mixed
tensors, does not claim Turing-completeness, and does not assign a new Zenodo
version DOI.
\end{remark}



\clearpage
\section{\QSCver{shifted.0.16} Release-Final DOI and Audit Hygiene}
\label{sec:v4016-final-local-audit-release-final-release-final-hygiene}

\begin{quote}\small
\noindent\textbf{Status.}
This \QSCver{shifted.0.16} layer is the final local audit and release-final hygiene
layer. It records the assigned version DOI 10.5281/zenodo.19811662, does not change any mathematical content,
does not modify the Historical Witness, and does not change any OP status. It
only checks and records the exact boundary between the frozen local release-final state
and a later DOI-bearing release-final pass.
\end{quote}

\subsection{Final local audit principle}
\label{subsec:v4016-final-local-audit-principle}

\begin{definition}[Final local audit principle]
\label{def:v4016-final-local-audit-principle}
A final local audit is a pre-release technical verification state. It verifies
source presence, PDF build status, audit-report availability, checksum
availability, release-final Zenodo description availability, and render-sample
availability before a DOI-bearing release-final pass is allowed.
\end{definition}

\begin{remark}[Release-final packaging record]
\label{rem:v4016-release-final-does-not-publish}
The release-final audit is not a Zenodo publication and is not a DOI-bearing release. It
is a local check that the next pass can be a release-final packaging pass after
a version DOI 10.5281/zenodo.19811662 is recorded.
\end{remark}

\subsection{Release-final checklist}
\label{subsec:v4016-release-final-release-final-checklist}

\begin{definition}[Release-final checklist]
\label{def:v4016-release-final-release-final-checklist}
The release-final audit is successful if the following checks hold:
\begin{enumerate}[leftmargin=*]
 \item Active Core source exists;
 \item Historical Witness source exists;
 \item Active Core PDF builds with zero fatal errors;
 \item Historical Witness PDF builds with zero fatal errors;
 \item undefined references, undefined citations, duplicate labels, overfull
 hbox warnings, and overfull vbox warnings are all zero;
 \item release-final Zenodo description exists;
 \item manifest and checksum files are generated;
 \item render contact sheet is generated;
 \item non-claim boundary remains explicit.
\end{enumerate}
\end{definition}

\subsection{Release-final packaging record}
\label{subsec:v4016-release-final-pass-input-contract}

\begin{definition}[Release-final packaging record]
\label{def:v4016-release-final-pass-input-contract}
The only missing input for a release-final pass after this release-final audit is a new
Zenodo version DOI. Once that DOI is supplied, the next pass should be a
metadata-only release-finalization pass unless the user explicitly requests
additional cleanup first.
\end{definition}

\begin{proposition}[Single-input release-final readiness]
\label{prop:v4016-single-input-release-final-readiness}
If the release-final checklist holds, then the local release-final state is release-final ready
up to the supplied Zenodo version DOI 10.5281/zenodo.19811662.
\end{proposition}

\begin{proof}
The release-final checklist verifies the required source, build, audit, packaging,
render, and non-claim conditions. The release-final input contract states that
the remaining missing datum is the assigned version DOI. Therefore, once that DOI
is supplied, release-final packaging can proceed without changing mathematical
content.
\end{proof}

\subsection{Final local non-claim boundary}
\label{subsec:v4016-final-local-non-claim-boundary}

\begin{definition}[Final local non-claim boundary]
\label{def:v4016-final-local-non-claim-boundary}
The final local audit preserves the following non-claims:
\[
\begin{array}{ll}
\text{OP~1/4} & \text{remains open},\\
\text{OP~5(B)} & \text{remains open},\\
\text{primitive mixed tensors} & \text{not activated},\\
\text{Turing-completeness} & \text{not claimed},\\
\text{standard physical calculations} & \text{not replaced},\\
\text{Zenodo version DOI} & \text{10.5281/zenodo.19811662}.
\end{array}
\]
\end{definition}

\subsection{Release-final audit conservativity}
\label{subsec:v4016-release-final-conservativity}

\begin{theorem}[Final local audit release-final conservativity]
\label{thm:v4016-final-local-audit-release-final-conservativity}
The \QSCver{shifted.0.16} release-final audit state is conservative over
the \QSCver{shifted.0.1}--\QSCver{shifted.0.15} cleanup sequence and over the public \QSCver{shifted.0.0} release-final
mathematical content.
\end{theorem}

\begin{proof}
The release-final audit adds audit and release-final readiness conditions only. It does
not modify $\Sigma_4$, $\mathrm B_4$, $\mathfrak D_4$, $\mathcal N_4$, tensor
grammar, readout compatibility, artifact roles, build rules, Historical Witness
content, or OP-status guards. Therefore it is conservative over the cleanup
sequence and over the public \QSCver{shifted.0.0} mathematical content.
\end{proof}

\begin{remark}[\QSCver{shifted.0.16} non-claim]
\label{rem:v4016-non-claim}
Version~4.0.16 is a release-final audit state. It
does not add new physics, does not close OP~1/4 or OP~5(B), does not activate
primitive mixed tensors, does not claim Turing-completeness, and records the assigned Zenodo version DOI 10.5281/zenodo.19811662.
\end{remark}



\clearpage
\section{\QSCver{shifted.0.16} Release-Final DOI and Packaging Record}
\label{sec:v4016-release-final-doi-packaging-record}

\begin{quote}\small
\noindent\textbf{Status.}
This section records the release-final DOI and packaging pass for
Version~4.0.16. The pass is metadata-only: it inserts the assigned Zenodo
version DOI, normalizes release-final artifacts, rebuilds the PDFs, regenerates
the audit records, and preserves all mathematical non-claims.
\end{quote}

\subsection{DOI record}
\label{subsec:v4016-release-final-doi-record}

\begin{definition}[Release-final DOI record]
\label{def:v4016-release-final-doi-record}
The release-final DOI record is:
\[
\begin{array}{ll}
\text{Concept DOI:} & 10.5281/zenodo.19210728,\\
\text{Prior public version DOI:} & 10.5281/zenodo.19805670,\\
\text{Version DOI:} & 10.5281/zenodo.19811662.
\end{array}
\]
\end{definition}

\begin{remark}[Version DOI is metadata]
\label{rem:v4016-version-doi-is-metadata}
The version DOI is a release and citation identifier. It is not a mathematical
operator, not a new object in the shifted-namespace normal form, and not an OP-closure
certificate.
\end{remark}

\subsection{Release-final conservativity}
\label{subsec:v4016-release-final-doi-pass-conservativity}

\begin{theorem}[Release-final DOI/pass conservativity]
\label{thm:v4016-release-final-doi-pass-conservativity}
The \QSCver{shifted.0.16} release-final DOI and packaging pass is conservative over the
\QSCver{shifted.0.16} local audit state.
\end{theorem}

\begin{proof}
The pass inserts the version DOI, normalizes release-final wording and file
names, rebuilds the PDFs, and regenerates audit and checksum artifacts. It does
not modify $\Sigma_4$, $\mathrm B_4$, $\mathfrak D_4$, $\mathcal N_4$, tensor
grammar, readout compatibility, artifact roles, build rules, Historical Witness
content, or OP-status guards. Therefore the pass is conservative.
\end{proof}

\begin{remark}[\QSCver{shifted.0.16} release-final non-claim]
\label{rem:v4016-release-final-non-claim}
Version~4.0.16 release-final does not add new physics, does not close OP~1/4
or OP~5(B), does not activate primitive mixed tensors, does not claim
Turing-completeness, and does not replace standard QFT/QED/Lifshitz or
general-relativistic calculations.
\end{remark}



\clearpage
